\documentclass[12pt]{article}
\usepackage[T1]{fontenc}
\usepackage[utf8]{inputenc}
\usepackage{lmodern}
\usepackage{textcomp}
\usepackage{enumitem}
\usepackage{geometry}
\usepackage{graphicx}
\usepackage{amsmath, amssymb}
\geometry{margin=1in}
\begin{document}
\begin{center}
History - K-12 Level Assessment 12 \
Total Marks: 40 \
Answer all questions.
\end{center}

\section*{Multiple Choice (10 marks)}
\begin{enumerate}[label=\arabic*.]
\item The 30 upright stones at Stonehenge are called \_\_\_\_\_\_\_; the horizontal stones are \_\_\_\_\_\_\_\_\_.
\begin{enumerate}[label=(\alph*)]
    \item legumes, sarsens
    \item lintels, stelae
    \item sarsens, lintels
    \item trilithon, tholoi
\end{enumerate}
\item A proposed explanation for some phenomenon that may be derived initially from empirical observation through a process called induction is a:
\begin{enumerate}[label=(\alph*)]
    \item deduction.
    \item theory.
    \item hypothesis.
    \item scientific method.
\end{enumerate}
\item Upon exploring Rhode Island, the Italian navigator Giovanni de Verrazzano noted that the native people resembled:
\begin{enumerate}[label=(\alph*)]
    \item Italians.
    \item Asians.
    \item Australians.
    \item Melanesians.
\end{enumerate}
\item In primates, the location of which of the following determines if it is a quadruped or biped?
\begin{enumerate}[label=(\alph*)]
    \item the feet
    \item the skull
    \item the foramen magnum
    \item the ulna
\end{enumerate}
\item What was the Chinese technique of stamped or pounded earth that was used in construction of house walls and defensive structures?
\begin{enumerate}[label=(\alph*)]
    \item Hang-t'u
    \item Scapulimancy
    \item Lung-shan
    \item Qin manufacture
\end{enumerate}
\end{enumerate}

\section*{True / False (10 marks)}
\textit{Answer True or False}
\begin{enumerate}[label=\arabic*.]
\setcounter{enumi}{5}
\item True or False: The levels of 13 C in bones can determine if someone ate mostly meat or vegetables.
\item True or False: The sagittal crest is a ridge of bone found on the forehead, not above the eye orbits.
\item True or False: Neandertals had smaller brains and larger noses than modern Homo sapiens, which is a common misconception.
\item True or False: The Adena culture was centered in the Ohio River valley around 2,800 years ago.
\item True or False: Diffusionism is the theory that characteristics of civilization spread from a central location to other regions.
\end{enumerate}

\section*{Long Form Response (20 marks)}
\textit{Answer each question with a clear and complete explanation.}
\begin{enumerate}[label=\arabic*.]
\setcounter{enumi}{10}
\item Discuss the factors that contributed to the preservation of the wooden spears found at the 300,000-year-old site in Schöningen, Germany, and their significance in understanding early human life.
\item Discuss the origins of Minoan civilization, focusing on the factors that contributed to its unique development on Crete, and analyze the influences of local and external cultures.
\end{enumerate}

\vfill
\noindent\textit{End of Paper}
\end{document}